% !TeX root = ../main.tex

\chapter{VHDL}

In this chapter the VHDL description of the Writing Master is presented, starting from an external view, with the entity description, to the internal architecture of the device.


\section{Entity declaration}

Listing \ref{list:entity_declaration} reports the external description of the Writing Master, as presented in Figure \ref{fig:i2c_wm_design}:

\lstinputlisting[firstline=13, lastline=24, caption={Writing Master Entity}, label={list:entity_declaration}]{../../src/writing_master.vhd}

The entity interfaces are presented by the \texttt{port} description, where all input and output signals are reported. As already said, the \texttt{sda} signal is a \texttt{inout} port because it is used both by the master and the slave. All the signals are based on the \texttt{std\_logic} family, and not \texttt{bit} type, because they can assume other logic values rather than only \texttt{1} or \texttt{0}, as the \texttt{Z} logic value for the \texttt{inout} signal. Every signal is represented by one bit, except for \texttt{addr} and \texttt{data} signals that are \texttt{std\_logic\_vector} because they are parallel buses on 7 and 8 bits respectively.

\section{Architecture description}

This section provides the internal description of the entity, including the declarative part and the statement part of the architecture.

\subsection{Declarative part}

Listing \ref{list:declarative_part} reports all the internal signals and types used to implement the logic of the Writing Master:

\lstinputlisting[firstline=26, lastline=35, caption={Architecture declarative part}, label={list:declarative_part}]{../../src/writing_master.vhd}

The type \texttt{state\_t} defines all the states of the FSM. The \texttt{curr\_state} and \texttt{next\_state} signals are declared as \texttt{state\_t} type to track the state. Regarding the counter signals, \texttt{scl\_count} and \texttt{bit\_count}, they are defined as \texttt{unsigned} type (5 and 4 bits respectively). Both signals are represented as a bus of wires, but they are interpreted as natural numbers allowing the use of basic arithmetic operations (\texttt{std\_logic\_vector} type is interpreted just as wires). Adding to this, the choice of the \texttt{unsigned} type, rather than \texttt{integer} type, is justified by the possibility to specify the exact number of bits needed (\texttt{integer} type occupies 32 bits) and to access the individual bits. The \texttt{addr\_signal} and \texttt{data\_signal} are used to register the values of the corresponding inputs when the \texttt{valid} signal is \texttt{1}. Finally, the \texttt{scl\_signal} and \texttt{sda\_signal} act as internal registers to update and hold the corresponding output signals.

\newpage

\subsection{Statement part}

The statement part is divided into three processes corresponding to the three blocks illustrated in Figure \ref{fig:block_diagram}:

\begin{itemize}
    \item \texttt{p\_STATE\_REG}, corresponding to the \texttt{REG\_STATE LOGIC} block;
    \item \texttt{p\_NEXT\_STATE\_LOGIC}, corresponding to the \texttt{NEXT\_STATE LOGIC} block;
    \item \texttt{p\_OUTPUT\_LOGIC}, corresponding to the \texttt{OUTPUT LOGIC} block.
\end{itemize}

These three processes are illustrated in detail in the following sections.


\subsubsection{p\_STATE\_REG}

\lstinputlisting[firstline=42, lastline=78, caption={p\_STATE\_REG process}, label={list:p_state_reg}]{../../src/writing_master.vhd}

